\tongjisetup{
  %******************************
  % 注意：
  %   1. 配置里面不要出现空行
  %   2. 不需要的配置信息可以删除
  %******************************
  % 秘级
  secretlevel={公开},
  secretyear={0},
  %
  %
  % 中文信息
  ctitle={基于雷视融合的吊装作业防碰撞监测 \\ 与安全状态评估方法研究}, % 可手动加入换行符 \\ 控制换行
  cauthor={陈宇凡},
  studentnumber={2330894},
  cdepartment={土木工程学院},
  csubjectcategory={工学}, % 学科门类
  cmajorfirst={土木水利}, % 专业学位类别
  cmajorsecond={土木工程}, % 专业领域
  cresearchfield={智能建造}, % 研究方向
  csupervisor={卢昱杰~教授},
  cassosupervisor={}, % 副指导老师/行业导师
  cjointinstitution={}, % 联合培养单位
  cdate={\zhdigits{2026}年\zhnumber{5}月}, % 手动指定日期, 否则自动使用当前时间
  cfunds={}, % 中文版基金信息
  % 英文信息
  etitle={Research on Anti-collision Monitoring and Safety State Assessment Method for Tower Crane Hoisting Based on LiDAR-Vision Fusion},
  eauthor={Yufan Chen},
  edepartment={College of Civil Engineering},
  esubjectcategory={Engineering}, % 学科门类
  emajorfirst={Civil and Hydraulic Engineering}, % 专业学位类别
  emajorsecond={Civil Engineering}, % 专业领域
  eresearchfield={Intelligent Construction}, % 研究方向
  esupervisor={Prof. Yujie Lu},
  eassosupervisor={}, % 副指导老师/行业导师
  ejointinstitution={}, % 联合培养单位
  edate={May, 2026}, % 手动指定日期, 否则自动使用当前时间
  efunds={} % 英文版基金信息
}

% 后处理
\makeatletter
\define@key{tongji}{cheadingtitlenew}{\cheadingtitlenew{#1}}
\gdef\cheadingtitlenew#1{\gdef\tongji@cheadingtitlenew{#1}}
\tongjideletelinebreakcmd{tongji@ctitle}{tongji@posttitle}
\cheadingtitlenew{\tongji@posttitle}
\makeatother

% 定义中英文摘要和关键字
\begin{cabstract}
  吊钩端近场空间是塔吊吊装过程中碰撞风险最集中的区域，实现其动态避障感知与实时预警对保障施工作业安全具有重要意义。现有塔臂端远距离观测方案受制于分辨率不足、遮挡频繁与近场盲区等固有局限，难以满足该区域的高可靠感知需求。将激光雷达与工业相机直接部署于吊钩端，可从物理层面缩短传感距离，提升近场感知精度与持续性。然而，吊钩端处于持续三维运动状态，其位姿随吊臂回转与钢丝绳伸缩实时变化，传感器观测基准随之动态改变。因此在这种高动态条件下，通过异构传感器融合实现环境与吊物状态的一致感知，并在此基础上建立有效的碰撞预警策略，是吊钩端感知亟待解决的关键技术难点。围绕该难点，本文的核心科学问题可归纳为：动态吊装场景中多源传感信息的时空一致性表征与碰撞风险协同判定方法。

  为此，本文构建了以双激光雷达与工业相机为核心的集成式吊钩端感知系统，围绕多传感器时空标定、融合点云定位与动态目标检测、雷视融合分级预警三个层面展开研究，形成了涵盖感知、定位、判级与控制建议的完整预警闭环。

  首先，完成双激光雷达与工业相机的一体化防护盒体设计与装配验证，采用双雷达侧向对称布置与相机垂直向下安装实现近场盲区互补与覆盖优化，通过PTPv2软件授时与融合层时间窗门控建立跨模态时间同步。采用粗配准初值构造与自动精配准优化相结合的两步法，实现双雷达外参标定，建立动态定位主链与静态感知支链相分离的分层式空间参考体系。其次，基于标定外参实现双雷达点云实时融合与坐标统一，通过扫描线拓扑重映射与时间窗口同步适配LIO-SAM紧耦合定位系统，输出连续位姿轨迹与静态地图。以配准点云为输入，经近场ROI裁剪后采用基于遮挡机理的M-detector进行逐点动态判别，结合欧式聚类与二维常速度卡尔曼滤波实现动态目标生成与短时跟踪。在视觉感知层面，采用YOLO26实例分割网络对吊装构件进行实时识别与像素级分割，通过TensorRT FP16半精度量化部署至Jetson Orin NX边缘平台。基于分割掩码将点云逐点投影至图像平面并回查实例点云集合，构建静态包络球与动态包络圆柱相结合的双基准安全包络模型。最后，建立基于三维净空的静态分级判据与基于当前净空、预测净空及碰撞时间（TTC）协同判定的动态风险升级机制，结合迟滞状态机输出结构化预警等级与避障控制建议。

  实际施工场景实验结果表明，点云投影关联精确率达92.17\%，F1分数84.69\%，验证了语义-几何跨模态映射的准确性；动态风险分级准确率达86.65\%，误报率5.35\%，漏报率8.00\%，在相向接近工况下可实现0.70\,s至1.60\,s的提前预警，表明基于净空与TTC的协同预警机制实现了可靠的分级判定。全链路端到端时延均值101.43\,ms，在典型吊装作业接近速度下具备充分的时效裕量。研究结果表明，本文所提方法能够在施工吊装场景下实现稳定的吊钩端近场感知与可靠的风险分级预警，为塔吊吊装过程的动态避障与安全预警提供了具有工程可行性的技术方案。
\end{cabstract}

\ckeywords{雷视融合，塔吊吊装，吊钩端感知，动态避障，风险预警}

\begin{eabstract}
  The near-field space around the hook end is the most collision-concentrated region during tower crane hoisting, and achieving dynamic obstacle avoidance perception and real-time early warning in this area is of great significance for construction safety. Existing long-range observation schemes deployed on the jib end suffer from inherent limitations such as insufficient resolution, frequent occlusion, and near-field blind zones, making it difficult to meet the high-reliability perception requirements of this region. Deploying LiDAR and industrial cameras directly at the hook end can physically shorten the sensing distance and improve near-field perception accuracy and persistence. However, the hook end undergoes continuous three-dimensional motion, with its pose changing in real time as the jib rotates and the wire rope extends and retracts, causing the sensor observation reference frame to change dynamically. How to achieve consistent representation of the environment and hoisted object states through heterogeneous sensor fusion under such highly dynamic conditions, and on this basis establish effective collision warning strategies, is a unique technical challenge that distinguishes hook-end perception from fixed-deployment schemes, and is also a key scientific problem in the field of hoisting safety.

  To address this problem, an integrated hook-end perception system comprising dual LiDARs and an industrial camera is developed, along with methods for spatiotemporal calibration, fused point cloud localization and dynamic object detection, and LiDAR-vision fusion-based graded warning, forming a complete closed loop of perception, localization, grading, and control suggestion.

  First, an integrated protective enclosure for the dual LiDARs and industrial camera is designed and assembly-verified, with the two LiDARs arranged symmetrically on the sides and the camera installed vertically downward to achieve near-field blind-spot complementarity and coverage optimization. Cross-modal temporal synchronization is established via PTPv2 software timestamping and a fusion-layer time window gating mechanism. A two-stage extrinsic calibration method combining coarse initial guess construction and automatic fine registration optimization is proposed for dual-LiDAR alignment, establishing a hierarchical spatial reference framework with separate dynamic localization main chain and static perception branch chain. Second, based on the calibrated extrinsic parameters, real-time dual-LiDAR point cloud fusion and coordinate unification are achieved; the fused point cloud is adapted to the LIO-SAM tightly-coupled localization system through scan-line topology remapping and time-window synchronization, producing continuous pose trajectories and a static map. Using the registered point cloud as input, the occlusion-based M-detector performs point-wise dynamic discrimination after near-field ROI cropping, and Euclidean clustering together with a 2D constant-velocity Kalman filter generates dynamic object candidates with short-term tracking. On the visual perception side, the YOLO26 instance segmentation network is deployed on a Jetson Orin NX edge platform via TensorRT FP16 half-precision quantization for real-time recognition and pixel-level segmentation of hoisting components. Based on the segmentation masks, point clouds are projected point-by-point onto the image plane to retrieve instance-level point cloud sets, and a dual-baseline safety envelope model combining a static bounding sphere and a dynamic bounding cylinder is constructed. Finally, a static clearance-based grading criterion and a dynamic risk escalation mechanism jointly determined by current clearance, predicted clearance, and time-to-collision (TTC) are established, with a hysteresis state machine outputting structured warning levels and obstacle avoidance control suggestions.

  Experimental validation in actual construction scenarios yields the following results: point cloud projection association achieves a precision of 92.17\,\% and an F1-score of 84.69\,\%, confirming the accuracy of semantic-to-geometric cross-modal mapping; dynamic risk grading accuracy reaches 86.65\,\% with a false-alarm rate of 5.35\,\% and a miss rate of 8.00\,\%, achieving early warning lead times of 0.70\,s to 1.60\,s in opposing approach scenarios, demonstrating that the collaborative warning mechanism based on clearance and TTC achieves reliable graded determination; the end-to-end pipeline latency averages 101.43\,ms, providing sufficient temporal margin at typical hoisting approach speeds. These results demonstrate that the proposed method can achieve stable hook-end near-field perception and reliable risk graded warning in construction hoisting scenarios, offering an engineering-feasible technical solution for dynamic obstacle avoidance and safety warning in tower crane hoisting operations.
\end{eabstract}

\ekeywords{LiDAR-vision fusion, tower crane hoisting, hook-end perception, dynamic obstacle avoidance, risk early warning}
